\section{Setting}\label{sec:setting}

We use units in which the magnetic field and the magnetic length are equal to one ($\hbar=e=1$).

\subsection{The torus and its lowest Landau level}

\begin{definition}[the torus and its lowest Landau level]\label{def:torus}
Let $N\in\N$ and $\Nphi:=2N$. Let $L_x,L_y>0$ with $L_xL_y=2\pi\Nphi=4\pi N$, put $\kappa:=2\pi/L_y$, so that
$L_x=\kappa\Nphi$, and let $\Tor_N:=\R^2/(L_x\Z\times L_y\Z)$, with points $r=(x,y)$, Lebesgue measure
$\dd r$, area $\abs{\Tor_N}=L_xL_y$, and flat distance $d(r,r')$. Let
\begin{equation}\label{eq:covariant}
D_x:=\partial_x,\qquad D_y:=\partial_y-\iu x,\qquad D:=\tfrac12(D_x-\iu D_y)
\end{equation}
be the covariant derivatives of the vector potential $(0,x)$, whose magnetic field is
$\partial_xA_y-\partial_yA_x=1$. The \emph{lowest Landau level} $\hN$ is the space of smooth functions
$f:\R^2\to\C$ such that
\begin{gather}
f(x,y+L_y)=f(x,y),\qquad f(x+L_x,y)=e^{\iu L_xy}f(x,y), \label{eq:bc}\\
(D_x+\iu D_y)f=0, \label{eq:lll}
\end{gather}
with the inner product $\ip{f}{g}:=\int_{[0,L_x)\times[0,L_y)}\overline{f}\,g\,\dd r$. By \eqref{eq:bc} the
functions $\abs{f}^2$ and $\overline fg$ are $L_x$- and $L_y$-periodic, so the integral is an integral over
$\Tor_N$.
\end{definition}

The conditions \eqref{eq:bc} are the magnetic periodicity conditions in this gauge without additional
Aharonov--Bohm phases through the two cycles of the torus; all results of this paper concern these
boundary conditions. The direction $x$ will be called \emph{transverse}: the Landau-gauge orbitals below are
localized in $x$ and extended in $y$.

\begin{lemma}[Landau-gauge orbitals]\label{lem:orbitals}
For $n\in\Z$ let $\phi_n(x,y):=(L_y\sqrt\pi)^{-1/2}e^{\iu\kappa ny}e^{-(x-\kappa n)^2/2}$, and for
$j\in\Z_{\Nphi}$ let
\begin{equation}\label{eq:psi}
\psi_j:=\sum_{m\in\Z}\phi_{\tilde\jmath+m\Nphi},
\end{equation}
where $\tilde\jmath\in\Z$ is any representative of $j$.
\begin{enumerate}[(i)]
\item The series \eqref{eq:psi} converges, together with all its derivatives, locally uniformly on $\R^2$;
$\psi_j$ does not depend on the choice of $\tilde\jmath$, and $\psi_j\in\hN$.
\item $(\psi_j)_{j\in\Z_{\Nphi}}$ is an orthonormal basis of $\hN$. In particular $\dim\hN=\Nphi$.
\item For $f\in\hN$, $Df=D_xf=\partial_xf$, and $Df$ satisfies \eqref{eq:bc}. In particular
$(D\phi_n)(x,y)=-(x-\kappa n)\,\phi_n(x,y)$.
\end{enumerate}
\end{lemma}

\begin{proof}
(i) Every derivative of $\phi_n$ is $\phi_n$ times a polynomial in $x$ and $n$, and
\[
\abs{\phi_n(x,y)}=(L_y\sqrt\pi)^{-1/2}e^{-(x-\kappa n)^2/2};
\]
these Gaussians are summable over $n\in\tilde\jmath+\Nphi\Z$ with any polynomial weight, locally uniformly in
$x$. This gives the convergence, and replacing $\tilde\jmath$ by $\tilde\jmath+\Nphi$ only relabels the sum. Since
$\kappa L_y=2\pi$, each $\phi_n$ is $L_y$-periodic in $y$. Since $\kappa\Nphi=L_x$,
\[
\phi_n(x+L_x,y)=(L_y\sqrt\pi)^{-1/2}e^{\iu\kappa ny}e^{-(x-\kappa(n-\Nphi))^2/2}=e^{\iu L_xy}\phi_{n-\Nphi}(x,y),
\]
and summing over $n\in\tilde\jmath+\Nphi\Z$ gives the second condition in \eqref{eq:bc} for $\psi_j$. Finally
$(D_x+\iu D_y)\phi_n=\bigl(-(x-\kappa n)+\iu(\iu\kappa n)+x\bigr)\phi_n=0$, and \eqref{eq:lll} follows
termwise.

(ii) Let $0\le i,j<\Nphi$ be representatives. For $n,n'\in\Z$ the $y$-integral of
$\overline{\phi_n}\phi_{n'}$ over $[0,L_y)$ is $L_y\delta_{nn'}(L_y\sqrt\pi)^{-1}e^{-(x-\kappa n)^2}$. Hence
$\ip{\psi_i}{\psi_j}=0$ for $i\ne j$, and
\[
\ip{\psi_j}{\psi_j}=\sum_{m\in\Z}\int_0^{L_x}\pi^{-1/2}e^{-(x-\kappa j-mL_x)^2}\,\dd x
=\pi^{-1/2}\int_\R e^{-u^2}\dd u=1 .
\]
For completeness, let $f\in\hN$. By the first condition in \eqref{eq:bc}, $f(x,y)=\sum_{n\in\Z}f_n(x)e^{\iu\kappa ny}$
with smooth $f_n$, the series converging with all derivatives. Condition \eqref{eq:lll} reads
$f_n'(x)+(x-\kappa n)f_n(x)=0$ for every $n$, so $f_n(x)=a_ne^{-(x-\kappa n)^2/2}$ with $a_n\in\C$. Comparing
the coefficients of $e^{\iu\kappa ny}$ in the second condition of \eqref{eq:bc} gives
$f_n(x+L_x)=f_{n-\Nphi}(x)$, that is $a_n=a_{n-\Nphi}$, because $x+L_x-\kappa n=x-\kappa(n-\Nphi)$. Hence
$f=(L_y\sqrt\pi)^{1/2}\sum_{j=0}^{\Nphi-1}a_j\psi_j$.

(iii) Condition \eqref{eq:lll} says $D_yf=\iu D_xf$, so $Df=\frac12(D_xf-\iu\cdot\iu D_xf)=D_xf$. The operator
$\partial_x$ commutes with translations and with multiplication by $e^{\iu L_xy}$, so $\partial_xf$ satisfies
\eqref{eq:bc}. The formula for $D\phi_n=\partial_x\phi_n$ is immediate.
\end{proof}

The orbital $\psi_j$ is concentrated near the line $x=\kappa j$ modulo $L_x$:
\[
\abs{\psi_j(x,y)}\le(L_y\sqrt\pi)^{-1/2}\sum_{m\in\Z}e^{-(x-\kappa\tilde\jmath-mL_x)^2/2}.
\]
Note that $Df$ is in general not in $\hN$; the derivative enters only through the evaluation of $Df$ at a
point.

\begin{lemma}[evaluation vectors]\label{lem:evaluation}
For $R=(R_x,R_y)\in\R^2$ let
\[
\ev R:=\sum_{j\in\Z_{\Nphi}}\overline{\psi_j(R)}\,\psi_j,\qquad
\evD R:=\sum_{j\in\Z_{\Nphi}}\overline{(D\psi_j)(R)}\,\psi_j\qquad\in\hN .
\]
\begin{enumerate}[(i)]
\item $\ip{\ev R}{f}=f(R)$ and $\ip{\evD R}{f}=(Df)(R)$ for every $f\in\hN$. Moreover
$\ev{R+(0,L_y)}=\ev R$ and $\ev{R+(L_x,0)}=e^{-\iu L_xR_y}\ev R$, and likewise for $\evD R$; in particular
$\lvert\ev R\rangle\langle\ev R\rvert$ and $\norm{\ev R}$, $\norm{\evD R}$ are functions of $R\in\Tor_N$.
\item $\int_{\Tor_N}\lvert\ev R\rangle\langle\ev R\rvert\,\dd R=\one_{\hN}$.
\item With $q_{pk}^2:=p^2L_x^2+k^2L_y^2$ for $(p,k)\in\Z^2$,
\begin{align*}
\norm{\ev R}^2&=\frac1{2\pi}\sum_{(p,k)\in\Z^2}e^{-q_{pk}^2/4}\cos(kL_yR_x+pL_xR_y),\\
\norm{\evD R}^2&=\frac1{4\pi}\sum_{(p,k)\in\Z^2}\Bigl(1-\frac{q_{pk}^2}2\Bigr)e^{-q_{pk}^2/4}\cos(kL_yR_x+pL_xR_y).
\end{align*}
\item Let $L_{\min}:=\min\{L_x,L_y\}\ge4$. Then for every $R$,
\[
\bigl\lvert2\pi\norm{\ev R}^2-1\bigr\rvert\le4.2\,e^{-L_{\min}^2/4}\le0.08,\qquad
\bigl\lvert4\pi\norm{\evD R}^2-1\bigr\rvert\le2.1\,L_{\min}^2e^{-L_{\min}^2/4}\le0.62 .
\]
\end{enumerate}
\end{lemma}

\begin{proof}
(i) Since $(\psi_j)$ is an orthonormal basis and the inner product is conjugate-linear in its first argument,
$\ip{\ev R}{f}=\sum_j\psi_j(R)\ip{\psi_j}{f}=f(R)$, and in the same way $\ip{\evD R}{f}=(Df)(R)$, because $D$
is linear. The behaviour under $R\mapsto R+(L_x,0)$ follows from \eqref{eq:bc}:
$\ip{\ev{R+(L_x,0)}}{f}=e^{\iu L_xR_y}f(R)=\ip{e^{-\iu L_xR_y}\ev R}{f}$ for every $f\in\hN$, and likewise for
$\evD R$, since $Df$ satisfies \eqref{eq:bc} (Lemma~\ref{lem:orbitals}(iii)). The $y$-periodicity is similar.

(ii) For $f,g\in\hN$, $\ip{f}{\int\lvert\ev R\rangle\langle\ev R\rvert g\,\dd R}=\int\ip{f}{\ev R}\ip{\ev R}{g}\dd R
=\int\overline{f(R)}g(R)\,\dd R=\ip fg$.

(iii) Write $R=(x,y)$. By \eqref{eq:psi},
$\norm{\ev R}^2=\sum_j\abs{\psi_j(R)}^2=\sum_{n\in\Z}\sum_{p\in\Z}\overline{\phi_n(R)}\phi_{n+p\Nphi}(R)$, the pairs
$(j,m)$ and $(j,m')$ being relabelled by $n=\tilde\jmath+m\Nphi$ and $p=m'-m$. Using $\kappa\Nphi=L_x$ and
$\tfrac12(u^2+(u-pL_x)^2)=(u-pL_x/2)^2+p^2L_x^2/4$ with $u=x-\kappa n$,
\[
\norm{\ev R}^2=\frac1{L_y\sqrt\pi}\sum_{p\in\Z}e^{\iu pL_xy}e^{-p^2L_x^2/4}\sum_{n\in\Z}e^{-(x-pL_x/2-\kappa n)^2}.
\]
By the Poisson summation formula, $\sum_ng(v-\kappa n)=\kappa^{-1}\sum_k\hat g(2\pi k/\kappa)e^{2\pi\iu kv/\kappa}$
with $\hat g(\xi)=\int g(u)e^{-\iu\xi u}\dd u=\sqrt\pi e^{-\xi^2/4}$ for $g(u)=e^{-u^2}$. Here
$2\pi/\kappa=L_y$, $L_y\kappa=2\pi$, and $e^{-\iu kL_y\cdot pL_x/2}=e^{-\iu\pi kp\Nphi}=1$ because $\Nphi$ is
even. Hence $\norm{\ev R}^2=(2\pi)^{-1}\sum_{(p,k)}e^{-q_{pk}^2/4}e^{\iu(kL_yx+pL_xy)}$, and the sum is real,
being symmetric under $(p,k)\mapsto(-p,-k)$. For $\evD R$ the same computation, with
$D\phi_n=-(x-\kappa n)\phi_n$, replaces $e^{-v^2}$ by $(v^2-p^2L_x^2/4)e^{-v^2}$, $v=x-pL_x/2-\kappa n$,
because $(x-\kappa n)(x-\kappa n-pL_x)=v^2-p^2L_x^2/4$. The Fourier transform of $v^2e^{-v^2}$ is
$\sqrt\pi(\tfrac12-\tfrac{\xi^2}4)e^{-\xi^2/4}$, which gives the second formula.

(iv) The terms with $(p,k)=(0,0)$ give $1$ in both formulas. For $(p,k)\ne(0,0)$ put $n:=p^2+k^2\ge1$, so that
$q_{pk}^2\ge nL_{\min}^2\ge16$. Since $u\mapsto e^{-u/4}$ and, for $u\ge4$, $u\mapsto\frac u2e^{-u/4}$ are
decreasing, and $\abs{1-u/2}\le u/2$ for $u\ge2$,
\begin{gather*}
\bigl\lvert2\pi\norm{\ev R}^2-1\bigr\rvert\le\sum_{(p,k)\ne0}e^{-nL_{\min}^2/4}
\le e^{-L_{\min}^2/4}\sum_{(p,k)\ne0}e^{-4(n-1)},\\
\bigl\lvert4\pi\norm{\evD R}^2-1\bigr\rvert\le\sum_{(p,k)\ne0}\frac{nL_{\min}^2}2e^{-nL_{\min}^2/4}
\le\frac{L_{\min}^2}2e^{-L_{\min}^2/4}\sum_{(p,k)\ne0}n\,e^{-4(n-1)},
\end{gather*}
using $e^{-(n-1)L_{\min}^2/4}\le e^{-4(n-1)}$. The first sum is $e^4\bigl((\sum_{p\in\Z}e^{-4p^2})^2-1\bigr)$, which is
less than $4.08$; the second is less than $4.15$ (its terms with
$n=1,2$ contribute $4+8e^{-4}$, and the rest less than $10^{-3}$). This gives the bounds; the numerical
values follow from $L_{\min}\ge4$.
\end{proof}

\subsection{Magnetic translations}

For a function $f$ on $\R^2$ let
\begin{equation}\label{eq:magtrans}
(\thx f)(x,y):=e^{\iu\kappa y}f(x-\kappa,y),\qquad(\thy f)(x,y):=f(x,y+L_y/\Nphi).
\end{equation}

\begin{lemma}[magnetic translations]\label{lem:magtrans}
\begin{enumerate}[(i)]
\item $\thx$ and $\thy$ restrict to unitary operators on $\hN$, with $\thx\psi_j=\psi_{j+1}$ and
$\thy\psi_j=e^{2\pi\iu j/\Nphi}\psi_j$.
\item On $\hN$, $\thy\thx=e^{2\pi\iu/\Nphi}\thx\thy$ and $\thx^{\Nphi}=\thy^{\Nphi}=\one$.
\item $\thx$ and $\thy$ commute with $D_x$ and $D_y$ on smooth functions, hence with $D$.
\item $\thx\ev R=e^{\iu\kappa R_y}\ev{R+\kappa\ex}$ and $\thx\evD R=e^{\iu\kappa R_y}\evD{R+\kappa\ex}$ for every
$R\in\R^2$, where $\ex=(1,0)$.
\end{enumerate}
\end{lemma}

\begin{proof}
(i) Directly from the definitions, $\thx\phi_n=\phi_{n+1}$ and $\thy\phi_n=e^{\iu\kappa nL_y/\Nphi}\phi_n=e^{2\pi\iu n/\Nphi}\phi_n$.
Summing over $n\in\tilde\jmath+\Nphi\Z$ gives the action on $\psi_j$ (the phase $e^{2\pi\iu n/\Nphi}$ depends only on
the class of $n$). Both operators therefore map the orthonormal basis $(\psi_j)$ onto an orthonormal basis up
to phases, so they restrict to unitary operators on $\hN$. (ii) follows from (i). (iii) For $\thy$ this is
clear, since $D_x$ and $D_y$ do not involve $y$ explicitly. For $\thx$, $D_x\thx f=\thx D_xf$, and
\begin{align*}
D_y(\thx f)(x,y)&=e^{\iu\kappa y}\bigl[(\partial_yf)(x-\kappa,y)+\iu\kappa f(x-\kappa,y)-\iu xf(x-\kappa,y)\bigr]\\
&=e^{\iu\kappa y}(D_yf)(x-\kappa,y)=(\thx D_yf)(x,y).
\end{align*}
(iv) The inverse of $\thx$ on $\hN$ is $(\thx^{-1}f)(x,y)=e^{-\iu\kappa y}f(x+\kappa,y)$, and $\thx^*=\thx^{-1}$ by (i).
For $f\in\hN$, $\ip{\thx\ev R}{f}=(\thx^{-1}f)(R)=e^{-\iu\kappa R_y}f(R+\kappa\ex)=\ip{e^{\iu\kappa R_y}\ev{R+\kappa\ex}}{f}$.
By (iii), $D\thx^{-1}f=\thx^{-1}Df$, and the same computation gives the formula for $\evD R$.
\end{proof}

Only $\thx$, the translation by one orbital across the cuts, enters the index theorem.

\subsection{The spinful Fock space and its sectors}

\begin{definition}[Fock space, sectors, the many-body translation]\label{def:fock}
Let $\kN:=\hN\otimes\C^2$, with the orthonormal basis $\psi_j\otimes\sigma$, $j\in\Z_{\Nphi}$,
$\sigma\in\{\up,\dn\}$, where $\up,\dn$ denote the standard basis of $\C^2$. Let $\Fock:=\Ff(\kN)$ be the
fermionic Fock space over $\kN$ (of dimension $4^{\Nphi}$), with annihilation operators $c(f)$, $f\in\kN$,
conjugate-linear in $f$ and satisfying the canonical anticommutation relations
$\{c(f),c(g)^*\}=\ip fg$, $\{c(f),c(g)\}=0$. Put
\[
c_{j\sigma}:=c(\psi_j\otimes\sigma),\qquad n_{j\sigma}:=c_{j\sigma}^*c_{j\sigma},\qquad
N_\sigma:=\sum_{j\in\Z_{\Nphi}}n_{j\sigma},\qquad S_z:=\tfrac12(\Nup-\Ndn).
\]
For $\nn=(n_\up,n_\dn)\in\{0,\dots,\Nphi\}^2$ the \emph{sector} $\sect\nn$ is the joint eigenspace
$\{\Psi:\Nup\Psi=n_\up\Psi,\ \Ndn\Psi=n_\dn\Psi\}$, and $\Psec\nn$ is the orthogonal projection onto it. An
operator $A$ on $\Fock$ is \emph{sector-preserving} if $[A,\Nup]=[A,\Ndn]=0$, and \emph{even} if it commutes
with the parity $(-1)^{\Nup+\Ndn}$. For a unitary $u$ on $\kN$, $\Ff(u)$ denotes its second quantization, the
unitary on $\Fock$ with $\Ff(u)c(f)\Ff(u)^*=c(uf)$ that fixes the vacuum. The \emph{many-body magnetic
translation} is $\Theta:=\Ff(\thx\otimes\one)^*$.
\end{definition}

\begin{lemma}[sectors and the many-body translation]\label{lem:sectors}
\begin{enumerate}[(i)]
\item $\Fock=\bigoplus_\nn\sect\nn$, an orthogonal sum, and
$\Psec\nn=\one_{\{n_\up\}}(\Nup)\,\one_{\{n_\dn\}}(\Ndn)$.
\item An operator $A$ is sector-preserving if and only if $A=\sum_\nn\Psec\nn A\Psec\nn$, if and only if
$[A,\Psec\nn]=0$ for every $\nn$. Sector-preserving operators form a unital $*$-subalgebra of the operators on
$\Fock$, and they are even. If $H$ is self-adjoint and sector-preserving, then $f(H)$ is sector-preserving for
every function $f$ on the spectrum of $H$, and for every $\nn$ and every Borel set $B\subset\R$ the spectral
projection of $H|_{\sect\nn}$ onto $B$, extended by zero to $\Fock$, equals $\one_B(H)\Psec\nn$.
\item $\Theta$ is unitary and sector-preserving, $\Theta^{\Nphi}=\one$, and for all $f\in\kN$, $j$ and $\sigma$,
\[
\Theta^*c(f)\Theta=c\bigl((\thx\otimes\one)f\bigr),\qquad\Theta^*c_{j\sigma}\Theta=c_{j+1,\sigma},\qquad
\Theta^*n_{j\sigma}\Theta=n_{j+1,\sigma}.
\]
\end{enumerate}
\end{lemma}

\begin{proof}
(i) $\Nup$ and $\Ndn$ are commuting self-adjoint operators with joint spectrum in $\{0,\dots,\Nphi\}^2$, and the
$\Psec\nn$ are their joint spectral projections. (ii) $A$ commutes with $\Nup$ and $\Ndn$ if and only if it
commutes with all their spectral projections, hence with every $\Psec\nn$; since $\sum_\nn\Psec\nn=\one$ this is
equivalent to $A=\sum_\nn\Psec\nn A\Psec\nn$. These conditions are stable under sums, products and adjoints.
A sector-preserving operator commutes with $\Nup+\Ndn$, hence with $(-1)^{\Nup+\Ndn}$. If $H$ is self-adjoint and
sector-preserving, then $H=\bigoplus_\nn H|_{\sect\nn}$, so $f(H)=\bigoplus_\nn f(H|_{\sect\nn})$; with
$f=\one_B$ this gives the last claim. (iii) $\Theta$ is unitary because $\thx\otimes\one$ is.
$\Theta^*c(f)\Theta=\Ff(\thx\otimes\one)c(f)\Ff(\thx\otimes\one)^*=c((\thx\otimes\one)f)$, and with
$f=\psi_j\otimes\sigma$ Lemma~\ref{lem:magtrans}(i) gives the other two formulas. Summing them over $j$ shows
$\Theta^*N_\sigma\Theta=N_\sigma$. Finally $\Theta^{\Nphi}=\Ff((\thx\otimes\one)^{\Nphi})^*=\one$ by
Lemma~\ref{lem:magtrans}(ii).
\end{proof}

The convention $\Theta^*n_{j\sigma}\Theta=n_{j+1,\sigma}$ is that of \cite[Section~5.2]{BBDF21} for the lattice
translation.

\subsection{The hollow-core Hamiltonian}

For $a,c\in\Z_{\Nphi}$ the function $\overline{\psi_a}\,\psi_c$ is smooth on $\Tor_N$ by \eqref{eq:bc}; let
$\nabla^2=\partial_x^2+\partial_y^2$.

\begin{definition}[the hollow-core Hamiltonian]\label{def:hamiltonian}
For $a,b,c,d\in\Z_{\Nphi}$ let
\begin{equation}\label{eq:Vabcd}
V_{abcd}:=4\pi\int_{\Tor_N}\overline{\psi_a}\psi_c\,\overline{\psi_b}\psi_d\,\dd r+4\pi\int_{\Tor_N}\overline{\psi_b}\psi_d\,\nabla^2\bigl(\overline{\psi_a}\psi_c\bigr)\,\dd r,
\end{equation}
and let
\begin{equation}\label{eq:HN}
H_N:=\frac12\sum_{\sigma,\sigma'\in\{\up,\dn\}}\ \sum_{a,b,c,d\in\Z_{\Nphi}}V_{abcd}\;
c_{a\sigma}^*c_{b\sigma'}^*c_{d\sigma'}c_{c\sigma}
\end{equation}
on $\Fock$. Equivalently, by Parseval's identity on $\Tor_N$,
\[
V_{abcd}=\frac1{\abs{\Tor_N}}\sum_{q\in\frac{2\pi}{L_x}\Z\times\frac{2\pi}{L_y}\Z}4\pi\bigl(1-\abs q^2\bigr)\,
\widehat{\overline{\psi_a}\psi_c}(-q)\,\widehat{\overline{\psi_b}\psi_d}(q),\qquad\hat f(q):=\int_{\Tor_N}f(r)e^{-\iu q\cdot r}\,\dd r ,
\]
that is, $H_N$ is the second quantization, restricted to the lowest Landau level, of the pair interaction
$\Vint=4\pi(\delta_{\Tor}+\nabla^2\delta_{\Tor})$, the periodization over $L_x\Z\times L_y\Z$ of the planar
interaction $V:=4\pi(\delta+\nabla^2\delta)$, whose Fourier transform is $4\pi(1-\abs q^2)$, $4\pi$ times the first
Laguerre polynomial evaluated at $\abs q^2$.
\end{definition}

The operator $H_N$ is self-adjoint, since $V_{abcd}=\overline{V_{cdab}}$, and it is sector-preserving, even,
and invariant under global spin rotations. It commutes with $\Theta$ (Lemma~\ref{lem:covariance}), and
$H_N=8\pi\int_{\Tor_N}\sum_\nu B_\nu(R)^*B_\nu(R)\,\dd R$ (Theorem~\ref{thm:localform}(a)); in particular
$H_N\ge0$.

The normalization of \eqref{eq:Vabcd} is the one in which a pair of electrons of relative angular momentum
one in the plane has energy one. To state this precisely, consider two particles in the planar lowest Landau
level for the same field, in the symmetric gauge: in complex coordinates $z_i=x_i+\iu y_i$, a two-particle
state is $\Psi(z_1,z_2)=G(z_1,z_2)e^{-(\abs{z_1}^2+\abs{z_2}^2)/4}$ with $G$ entire and $\Psi$
square-integrable. With $Z:=(z_1+z_2)/2$ and $v:=z_1-z_2$ one has $\abs{z_1}^2+\abs{z_2}^2=2\abs Z^2+\abs v^2/2$,
and we expand $G(Z+v/2,Z-v/2)=\sum_{m\ge0}G_m(Z)v^m$. The component of relative angular momentum $m$ is
$\Psi_m:=G_m(Z)\,v^m\,e^{-\abs Z^2/2-\abs v^2/8}$; the $\Psi_m$ are mutually orthogonal and $\Psi=\sum_m\Psi_m$.
The quadratic form of $V$ on such states is
\[
\mathsf V(\Psi):=4\pi\int_{\R^2}\Bigl(F_\Psi(R,0)+(\nabla_r^2F_\Psi)(R,0)\Bigr)\dd R,\qquad
F_\Psi(R,r):=\abs{\Psi(R+r/2,R-r/2)}^2,
\]
the expectation of $V(r_1-r_2)$ in the sense of distributions; it depends only on $\abs\Psi^2$ and so is the
same in every gauge.

\begin{lemma}[$V_1$ units]\label{lem:units}
For every planar two-particle lowest-Landau-level state $\Psi$, $\mathsf V(\Psi)=\norm{\Psi_1}^2$. That is, on the
planar two-particle lowest Landau level the interaction $V$ is the orthogonal projection onto relative
angular momentum one: its pseudopotentials are $V_m=\delta_{m1}$.
\end{lemma}

\begin{proof}
In the coordinates $R=Z$ and $r=v$ (the change of variables $(r_1,r_2)\mapsto(R,r)$ has Jacobian one),
$F_\Psi(Z,v)=e^{-\abs Z^2}e^{-\abs v^2/4}\abs{\sum_mG_m(Z)v^m}^2$. At $v=0$ this is $e^{-\abs Z^2}\abs{G_0(Z)}^2$.
With $\nabla_v^2=4\partial_v\partial_{\bar v}$ and
$e^{-v\bar v/4}\abs{\sum_mG_mv^m}^2=\abs{G_0}^2+G_0\overline{G_1}\bar v+\overline{G_0}G_1v
+\bigl(\abs{G_1}^2-\tfrac14\abs{G_0}^2\bigr)v\bar v+(\text{terms in }v^2,\bar v^2,\text{ and of order }\abs v^3)$,
one finds $(\nabla_v^2F_\Psi)(Z,0)=e^{-\abs Z^2}\bigl(4\abs{G_1(Z)}^2-\abs{G_0(Z)}^2\bigr)$. Hence
\[
\mathsf V(\Psi)=16\pi\int_{\R^2}\abs{G_1(Z)}^2e^{-\abs Z^2}\dd Z
=\int\abs{G_1(Z)}^2e^{-\abs Z^2}\dd Z\int_{\R^2}\abs v^2e^{-\abs v^2/4}\dd v=\norm{\Psi_1}^2 .\qedhere
\]
\end{proof}

\begin{remark}[orbital form and normalization]\label{rem:orbitalform}
Theorem~\ref{thm:localform}(c) below writes $H_N$ in the orbital basis:
$H_N=(\kappa^3/\sqrt{2\pi})\sum_{n\in\Z_{2\Nphi}}\sum_\nu(\bcen\nu n)^*\bcen\nu n$, with the pair operators
$\bcen\nu n$ of Definition~\ref{def:pairs}. The operator $\sum_n\sum_\nu(\bcen\nu n)^*\bcen\nu n$ differs from
$H_N$ by the factor $\sqrt{2\pi}/\kappa^3$, which grows like $N^{3/2}$ at fixed aspect ratio. Every statement
of this paper about gaps that are uniform in $N$ refers to $H_N$ in the normalization of
Definition~\ref{def:hamiltonian}.
\end{remark}

\subsection{The zero modes}

The kernel of $H_N$ on the states with $N$ electrons is known exactly. The proof of the following theorem is given in
Appendix~\ref{app:kernel}.

\begin{theorem}[the zero modes]\label{thm:kernel}
Let $N\in\N$ and $L_x,L_y>0$ with $L_xL_y=4\pi N$, and let $\mathcal K_N$ be the kernel of $H_N$ restricted to
the states with $N$ electrons, $\bigoplus_{n_\up+n_\dn=N}\sect{(n_\up,n_\dn)}$.
\begin{enumerate}[(i)]
\item $\mathcal K_N\cap\sect{(n_\up,n_\dn)}=\{0\}$ whenever $\abs{n_\up-n_\dn}\ge2$.
\item If $N$ is odd, $\mathcal K_N\cap\sect{((N+1)/2,(N-1)/2)}$ and $\mathcal K_N\cap\sect{((N-1)/2,(N+1)/2)}$ are
two-dimensional.
\item If $N$ is even, $\mathcal K_N\subset\sect{(N/2,N/2)}$, $\dim\mathcal K_N=10$, and $\mathcal K_N$ consists of
spin singlets.
\end{enumerate}
\end{theorem}

Appendix~\ref{app:kernel} proves the theorem for tori of arbitrary modular parameter; we use it only for rectangular
tori. The kernel in (iii) is spanned by the torus states of Read and Rezayi \cite[Section~V]{RR96}, and the
kernel in (ii) by the corresponding states with one unpaired electron; we use only the dimensions. Since
$H_N\ge0$, the eigenvalues $E_1(N)\le E_2(N)\le\cdots$ of $H_N$ in the sector $(N/2,N/2)$ at even $N$,
repeated according to multiplicity, satisfy
\begin{equation}\label{eq:E11}
E_1(N)=\cdots=E_{10}(N)=0<E_{11}(N),
\end{equation}
and at odd $N$ the eigenvalues in each of the sectors $((N\pm1)/2,(N\mp1)/2)$ satisfy $E_1(N)=E_2(N)=0<E_3(N)$.
(The sector $(N/2,N/2)$ has dimension $\binom{2N}{N/2}^2\ge16$ for even $N\ge2$.)

\begin{remark}[the Hamiltonian of Appendix~\ref{app:kernel}]\label{rem:kernel-identification}
In its own notation, Appendix~\ref{app:kernel} defines the hollow-core Hamiltonian as $\sum_{i<j}\mathcal P_{ij}$, where $\mathcal P$ is the sum over
the lattice of periods of the translates of the planar projector onto relative angular momentum one, on the torus
$\C/(L_1\Z+L_1\tau\Z)$, with coordinates $X+\iu Y$, in the Landau gauge $(-Y,0)$. For $L_1=L_y$ and $\tau=\iu L_x/L_y$ this is our torus turned by a
quarter turn. The beginning of Appendix~\ref{app:kernel} shows, using Theorem~\ref{thm:localform}(c), that the
corresponding unitary identification of the two lowest Landau levels, which takes the orbital $\psi_j$ to the $j$-th
Landau-gauge orbital of Appendix~\ref{app:kernel}, carries $H_N$, with the spin unchanged, to that operator. The
identification preserves $\Nup$ and $\Ndn$, so Theorem~\ref{thm:kernel} holds for $H_N$ as stated.
\end{remark}
